
\section{标签中的单位}
\label{sec:units}

\begingroup
\def\pgfplotsmanualcurlibrary{units}

{\noindent {\emph{作者：Nick Papior Andersen}}}

\begin{pgfplotslibrary}{units}
        \index{Tick Labels!Units (km, m, etc)}%
    本库用于在标签中自动排版单位。本库借助不同的键，以一致的方式排版最终输出。
    调用其中任一命令会自动设置键 `|use units=true|'，因此无需额外操心。
\end{pgfplotslibrary}

\PGFPlots{} 具备支持单位的能力。这既便于快速自定义绘图，也便于在标签中加入单位。

加载本库会自动启用标签中的单位排版。目前它仅支持预定义的 SI 单位，但也实现了
用户自定义，因此可以按你喜欢的任何方式使用。

首先是用于开启或关闭单位系统的键。

\begin{pgfplotskey}{use units=\marg{boolean} (initially true)}
    该键用于让 \PGFPlots{} 启用下文所述的功能。加载本库时，该键会被置为 true。
    你可以用它临时决定是否在绘图中使用单位库。
\end{pgfplotskey}

\begin{pgfplotsxykey}{\x\ unit=\marg{unit} (initially empty)}
    这些键用于在对应坐标轴中设置单位。例如在 SI 单位制下，你可以将
    |x unit| 设为牛顿，即 |x unit=N|。
\end{pgfplotsxykey}

\begin{pgfplotsxykey}{\x\ unit prefix=\marg{prefix} (initially empty)}
    这些键用于设置单位的前缀。若 |y axis| 上的数值以千计，则应设置
    |y unit prefix=k|。前缀将被排版在单位之前。

    该命令不会干预坐标轴系统的基数。也就是说，刚刚提到的前缀不会把每个
    |y axis| 数值除以 1000。若需要这样做，请参见键 \meta{axis}| SI prefix|，
    见第~\ref{sec:SI:prefix} 节。

    注意：若未设置 \meta{axis}| unit|，则整个单位都不会被排版。


    \paragraph{注意事项：}

    请记住，标签的所有排版都发生在数学模式内（即在 \$ 定界符之内）。因此可以使用
    \textbackslash|frac| 以及其他数学命令。
\end{pgfplotsxykey}

通常只需使用上述键即可。它们是 \PGFPlots{} 提供的单位排版系统的基础。
%
% \usepgfplotslibrary{units}
\begin{codeexample}[]
\begin{tikzpicture}
\begin{axis}[use units,
    x unit=m,x unit prefix=k,
    y unit=N,y unit prefix=m,
    xlabel=Distance,ylabel=Force,
]
    \addplot coordinates {
        (1,2.3)
        (2,2.7)
        (3,2.1)
        (4,1.8)
        (5,1.5)
        (6,1.1)
    };
\end{axis}
\end{tikzpicture}
\end{codeexample}

下面的例子展示了按如下样式设置后得到的效果：
%
\begin{codeexample}[code only]
% x label becomes ``Temperature [T]'', y label becomes ``Nothing''
\pgfplotsset{use units,x unit=T,xlabel=Temperature,ylabel=Nothing}
% x label becomes ``Temperature'', y label becomes ``Nothing''
\pgfplotsset{use units,x unit prefix=m,xlabel=Temperature,ylabel=Nothing}
\end{codeexample}
%
请注意第二个例子。仅设置前缀并不会启用单位排版。因此，若希望排版标签，
应确保使用 |x unit| 键。

排版单位时，也可以修改其外观。例如，你可能不喜欢包围单位的方括号。幸运的是，
这些都可以通过下面的键来修改。

\begin{pgfplotskeylist}{unit marking pre=\marg{pre} (initially \textbackslash left[),unit marking post=\marg{post} (initially \textbackslash right]),unit markings=\mchoice{parenthesis,square brackets,slash space} (initially square brackets)}
    这些键用于设置单位的包围符号。初始设置得到
    $\left[\frac{1}{2}\right]$，以便在单位中排版分式。请注意，只有使用
    \textbackslash|dfrac| 才能得到较大的分式。可以通过选项键 |unit markings|
    轻松设置，其选项排版效果如下：
    %
\begin{codeexample}[code only]
\pgfplotsset{x unit=T,unit markings=parenthesis} % x unit becomes `` \left(T\right)''
\pgfplotsset{x unit=T,unit markings=square brackets} % x unit becomes `` \left[T\right]''
\pgfplotsset{x unit=T,unit markings=slash space} % x unit becomes `` / T''
\end{codeexample}
    %
    注意，所有单位排版都会先插入一个空格，然后再插入 |unit marking pre| 代码。

    当然，你也可以通过 |unit marking pre| 与 |unit marking post| 键分别手动设置。
    只需记住它们都在 \$ 定界符内排版。
\end{pgfplotskeylist}

通常会用特定字体排版单位。为此，本库提供了一个用于修改排版命令的选项。

\begin{pgfplotscodetwokey}{unit code}%\textbackslash mathrm\{\})}
    该键可被广泛使用。默认情况下，单位排版为
    |\mathrm{|\meta{unit prefix}\meta{unit}|}|。但如果你希望利用 |siunitx| 宏包
    （它在排版单位、数字与角度方面能力强大），只需将该键设置为
    %
\begin{codeexample}[code only]
\pgfplotsset{unit code/.code 2 args={\si{#1#2}}}
\end{codeexample}
    %
    这将使单位排版为 |\si{|\meta{unit prefix}\meta{unit}|}|。

    第一个参数排版为 \meta{unit prefix}，第二个参数为 \meta{unit}。

    最重要的是，该命令需要恰好两个参数。因此，如果你希望用粗体排版前缀、用
    常规罗马字体排版单位，则应调用
    %
\begin{codeexample}[code only]
\pgfplotsset{unit code/.code 2 args={\mathbf{#1}\mathrm{#2}}
\end{codeexample}
\end{pgfplotscodetwokey}


\subsection{预设 SI 前缀}
\label{sec:SI:prefix}
    \index{prefixes!unit}

为支持 SI 单位制，本库定义了若干预设键。这应能让前缀的指定更加直观，并提供更多
功能。例如，它提供了一种简便的缩放机制。

\begin{pgfplotsxykeylist}{%
    \x\ SI prefix=\mchoice{
        yocto,\dots,milli,centi,deci,deca,hecto,kilo,\dots,yotta%
    } (initially none),
    change \x\ base=\mchoice{true,false} (initially false)%
}
    这些键用于设置单位的前缀。允许的前缀为：
    %
    \begin{center}
        \begin{tabular}{>{\ttfamily}c>{$}c<{$}}
                \toprule
            \textrm{前缀} & $幂次$ \\
                \midrule
            yocto           & -24     \\
            zepto           & -21     \\
            atto            & -18     \\
            femto           & -15     \\
            pico            & -12     \\
            nano            & -9      \\
            micro           & -6      \\
            milli           & -3      \\
            centi           & -2      \\
            deci            & -1      \\
                \bottomrule
        \end{tabular}
            \qquad\qquad
        \begin{tabular}{>{\ttfamily}cc}
                \toprule
            \textrm{前缀} & 幂次 \\
                \midrule
            deca             & 1     \\
            hecto            & 2     \\
            kilo             & 3     \\
            mega             & 6     \\
            giga             & 9     \\
            tera             & 12    \\
            peta             & 15    \\
            exa              & 18    \\
            zetta            & 21    \\
            yotta            & 24    \\
                \bottomrule
        \end{tabular}
    \end{center}
    %
    此外，若设置键 |change |\meta{axis}| base=true|，还会重置坐标轴的基数。
    请务必\textbf{先}设置 |change |\meta{axis}| base|，\textbf{再}使用
    \meta{axis}| SI prefix| 键。

    用法见下面的示例。
    %
% \usepgfplotslibrary{units}
\begin{codeexample}[]
\begin{tikzpicture}
\begin{axis}[
    change x base,
    x SI prefix=kilo,x unit=m,
    y SI prefix=milli,y unit=N,
    xlabel=Distance,ylabel=Force,
]
    \addplot coordinates {
        (1000,1)
        (2000,1.1)
        (3000,1.2)
        (4000,1.3)
    };
\end{axis}
\end{tikzpicture}
\end{codeexample}
    %
    注意，|x axis| 已改变基数，但未显示 $\cdot 10^{3}$。这是通过键
    |change x base| 实现的。即使你使用了键 |y SI prefix=milli|，|y axis| 的基数
    也并未改变。请尝试在 |change x base| 之后紧接着加上 |change y base|，看看
    结果如何！
\end{pgfplotsxykeylist}

上述键是基数切换的简易实现。下面是对基数切换的进一步自定义，它便于实现带有
自定义基数切换的前缀。

\begin{pgfplotskey}{axis base prefix=%
    {axis \marg{axis} base \marg{base} prefix \marg{prefix}} (initially empty)%
}
    可以利用该键进一步自定义基数与前缀设置。
    %
\begin{codeexample}[code only]
\pgfplotsset{change x base,axis base prefix={axis x base -3 prefix k}}
\pgfplotsset{change x base,x SI prefix=kilo}
\end{codeexample}
    %
    因此上述两条命令是等价的。请记住，基数的方向应与前缀相反！
\end{pgfplotskey}
\endgroup
